% from ega0-20-fr.tex
\oldpage[0\textsubscript{IV-1}]{153}
complet en vertu de l'hypothèse, donc $E=D_C(L)$, muni de la topologie produit,
est métrisable et complet ; il est en outre admissible, car si $\mathfrak{R}$
est un idéal de définition dans $C$, la suite des
$(\mathfrak{R}\oplus L)^n=\mathfrak{R}^n\oplus\mathfrak{R}^{n-1}L$ tend vers
$0$. Comme l'application composée
\[
  D:B\xrightarrow{d_{B/A}}\Omega^1_{B/A}\to
  \widehat{\Omega}^1_{B/A}\widehat\otimes_{\widehat B}\widehat C=L
\]
est une $A$-dérivation continue de $B$ dans $L$, le $A$-homomorphisme continu
\[
  f:x\longmapsto(v(x),D(x))
\]
de $B$ dans $E$ définit sur $E$ une structure de $B$-extension. Comme $L$ est
un idéal fermé dans $E$, il résulte de \sref{0.19.9.5-fr} et de l'hypothèse
que l'application identique $C\to E/L$ (qui est un $B$-homomorphisme) se
factorise en $C\xrightarrow{g}E\to E/L$, où $g$ est un homomorphisme continu
tel que $g\circ v=f$ ; par suite \sref{0.20.1.3-fr}, $g$ est de la forme
$y\mapsto(y,D'(y))$, où $D'$ est une $B$-dérivation continue de $C$ dans $L$,
autrement dit $D'\circ v=D$. Compte tenu de \sref{0.20.7.14.4-fr}, on a
$D'=h\circ\widehat d_{C/A}$, où
$h:\widehat{\Omega}^1_{C/A}\to L$ est un $C$-homomorphisme continu ; mais on
a $\widehat d_{C/A}\circ v=\widehat v_{C/B/A}\circ D$ par définition, et
comme l'image de $B$ par $D$ engendre (topologiquement) le $C$-module $L$
\sref{0.20.4.7-fr}, la relation $D'\circ v=D$ donne bien
$h\circ\widehat v_{C/B/A}=1_L$.
\end{proof}

\begin{corollary}[20.7.19]
\label{0.20.7.19-fr}
Sous les hypothèses de \sref{0.20.7.18-fr}, si l'on suppose en outre que dans
$B$ et $C$ le carré d'un idéal ouvert soit ouvert, alors, pour que $C$ soit
une $B$-algèbre formellement lisse relativement à $A$, il faut et il suffit
que $\widehat v_{C/B/A}$ soit inversible à gauche.
\end{corollary}

\begin{proof}
En effet, les topologies de
$\widehat{\Omega}^1_{B/A}\widehat\otimes_{\widehat B}\widehat C$ et de
$\widehat{\Omega}^1_{C/A}$ sont alors déduites de celle de $C$
\sref{0.20.4.5-fr}, et tout $C$-homomorphisme de l'un dans l'autre est
nécessairement continu.
\end{proof}

\begin{env}[20.7.20]
\label{0.20.7.20-fr}
Soient $A$ un anneau topologique, $B$ une $A$-algèbre topologique
\emph{métrisable et complète}, $\mathfrak{R}$ un idéal \emph{fermé} de $B$,
$C=B/\mathfrak{R}$ l'anneau topologique quotient, qui est métrisable et
complet. Soit $(\mathfrak{M}_k)$ un système fondamental décroissant de
voisinages de $0$ dans $B$ formé d'idéaux, et posons
$B_k=B/\mathfrak{M}_k$,
$\mathfrak{R}_k=(\mathfrak{R}+\mathfrak{M}_k)/\mathfrak{M}_k$,
$C_k=B_k/\mathfrak{R}_k$. On a un système projectif d'homomorphismes
\[
  \delta_{C_k/B_k/A}:\mathfrak{R}_k/\mathfrak{R}_k^2\to
  \Omega^1_{B_k/A}\otimes_{B_k}C_k
\]
\sref{0.20.5.11.3-fr}, d'où l'on déduit en passant à la limite un
homomorphisme canonique
\[
  \widehat\delta_{C/B/A}:\mathfrak{R}/\mathfrak{R}^2\to
  \widehat{\Omega}^1_{B/A}\widehat\otimes_{\widehat B}\widehat C
\]
et en raisonnant comme dans \sref{0.20.7.17-fr}, on voit que l'homomorphisme
canonique
$\widehat u_{C/B/A}:\widehat{\Omega}^1_{B/A}
\widehat\otimes_{\widehat B}\widehat C\to\widehat{\Omega}^1_{C/A}$ est
\emph{surjectif} et que $\operatorname{Im}(\widehat\delta_{C/B/A})$ est
\emph{dense} dans $\operatorname{Ker}(\widehat u_{C/B/A})$.
\end{env}

% (continued in ega0-21-fr.tex, before its first page marker)
\section{Différentielles dans les anneaux de caractéristique p}
Les résultats du présent paragraphe, de nature plus spéciale et technique que
ceux des §§19, 20 et 22, ne serviront qu'exceptionnellement dans le cours du
chap. IV. Leur rôle principal est ici dans la démonstration de trois théorèmes
du §22 (22.3.3, 22.5.8 et 22.7.3), dont le premier et le dernier interviennent
de façon essentielle dans la théorie «fine» des anneaux locaux noethériens du
chap. IV, §7.
